\documentclass[11pt]{amsart}
\usepackage{calc,amssymb,amsmath,ulem,stmaryrd,fullpage}
\usepackage{enumerate}
\usepackage{alltt}
\RequirePackage[dvipsnames,usenames]{color}
\let\mffont=\sf
\normalem
%\input{kmacros3.sty}
\input{xy}
\xyoption{all}
\usepackage{hyperref}
\usepackage{graphicx}
\usepackage[all,cmtip]{xy}
\usepackage{verbatim}
\usepackage[left=0.95in,top=0.95in,right=0.95in,bottom=0.95in]{geometry}
\newcommand{\tauCohomology}{T}
\newcommand{\FCohomology}{S}


\theoremstyle{theorem}
%\newtheorem*{mainthm*}{Main Theorem}

\newcommand{\note}[1]{\marginpar{\sffamily\tiny #1}}

\def\todo#1{\textcolor{Mahogany}%
{\sffamily\footnotesize\newline{\color{Mahogany}\fbox{\parbox{\textwidth-15pt}{#1}}}\newline}}

\renewcommand{\O}{\mathcal O}
\newcommand{\ie}{{\itshape i.e.} }
\newcommand{\roundup}[1]{\lceil #1 \rceil}
\newcommand{\rounddown}[1]{\lfloor #1 \rfloor}
\renewcommand{\to}[1][]{\xrightarrow{\ #1\ }}
\newcommand{\onto}[1][]{\protect{\xrightarrow{\ #1\ }\hspace{-0.8em}\rightarrow}}
\newcommand{\into}[1][]{\lhook \joinrel \xrightarrow{\ #1\ }}
%\newcommand{DBDual}[1]{{\underline{\omega}_{#1}}}
\newcommand{\DBDual}[1]{{{\uuline \omega}{}^{\mydot}_{#1}}}
\newcommand{\Tt}{{\mathfrak{T}}}
%\newcommand{\bn}{{\mathfrak{n}} }
\newcommand{\sol}[1]{ \vskip 6pt{\color{blue} {\bf Solution:} #1} \vskip 6pt}

\newcommand{\bR}{{\mathbb{R}}}
\newcommand{\va}{{\vec{a}}}
\newcommand{\vb}{{\vec{b}}}
\newcommand{\vzero}{{\vec{0}}}
\newcommand{\vi}{{\vec{i}}}
\newcommand{\vj}{{\vec{j}}}
\newcommand{\vk}{{\vec{k}}}
\newcommand{\vh}{{\vec{h}}}
\newcommand{\vr}{{\vec{r}}}
\newcommand{\vu}{{\vec{u}}}
\newcommand{\vv}{{\vec{v}}}
\newcommand{\vw}{{\vec{w}}}
\newcommand{\vx}{{\vec{x}}}
\newcommand{\vy}{{\vec{y}}}
\newcommand{\vz}{{\vec{z}}}
\newcommand{\vT}{{\vec{T}}}
\newcommand{\vN}{{\vec{N}}}
\newcommand{\vF}{{\vec{F}}}
\newcommand{\vG}{{\vec{G}}}
\newcommand{\bF}{\mathbb{F}}
\newcommand{\bQ}{\mathbb{Q}}
\newcommand{\bZ}{\mathbb{Z}}
\newcommand{\bC}{\mathbb{C}}
\newcommand{\Gal}{\text{Gal}}


\begin{document}
\title{HW \#2 -- Math 6320\\Spring 2015}
\author{Due: Thursday February 5th}
\maketitle

\begin{enumerate}
\item Let $E = \bQ(2^{1/2}, 2^{1/3})$.  Find all subfields of $E$.  
\item Same setup as in the previous problem.  Show that $E = \bQ(2^{1/2} + 2^{1/3})$.  
\item Let $K=\mathbb{Q}\left(\frac{-1+\sqrt{-3}}{2}\right)$. Give an example of two non-isomorphic field extensions $L_1$ and $L_2$
of $K$ such that $\mathrm{Gal}(L_1/K)\cong \mathrm{Gal}(L_2/ K)\cong \mathbb{Z}/3\mathbb{Z}$. Justify your claims. 
\item  Suppose that $f \in \bQ[x]$ is an irreducible polynomial and that $\alpha,\beta \in \bC$ are roots of $f$.  Suppose that $\bQ \subseteq K \subseteq \bC$ is such that $K/\bQ$ is a finite Galois extension.  Show that $\bQ[\alpha] \cap K$ is isomorphic to $\bQ[\beta] \cap K$.
\vskip 3pt
\emph{Hint:} We know there is an isomorphism $\sigma : \bQ[\alpha] \to \bQ[\beta]$ sending $\alpha$ to $\beta$.  Show that $\sigma$ extends to an automorphism of some larger field that sends $K$ to $K$.  
\item Suppose that $\alpha \in  \bC$ with $\alpha^n \in \bQ$ such that $\bQ[\alpha] / \bQ$ is Galois. Further
suppose that $F$ is the field containing $\bQ$ generated by all the roots of unity
in $\bQ[\alpha]$. Show that $\Gal(\bQ[\alpha] : F )$ is a cyclic group.
\item Suppose that $K, L \subseteq \bC$ are subfields, each being Galois over $\bQ$. Show that the
field $KL$, generated by $K$ and $L$, is Galois over $\bQ$. 
\item Same setup as in the previous problem.  Now suppose that $[K : \bQ]$ and $[L : \bQ]$ are relatively prime.  Compute $\Gal(KL/\bQ)$ in terms of $\Gal(K/\bQ)$ and $\Gal(L/\bQ)$.  Deduce that there is a field $E$ Galois over $\bQ$ with $[E : \bQ] = 44$.
\end{enumerate}

\end{document}

